\documentclass{article}
\usepackage{hyperref}
\usepackage{Style}

\nocite{*} % Comentar si quiero citar
%\addbibresource{bibliografia.bib} % Quitar el comentado si quiero usar bibliografia

\begin{document}

\begin{minipage}{2cm}
    \includegraphics[width=2cm]{imagen_puc.jpg}
\end{minipage}
\begin{minipage}{14cm}
    {\sc Pontificia Universidad Católica de Chile\\
    Facultad de Matemáticas\\
    Departamento de Matemática\\
    Profesor: Manuel Cabezas -- Ayudante: Benjamín Mateluna}
\end{minipage}
\vspace{1ex}

{\centerline{\bf Introducción a la Geometría - MAT1304}
\centerline{\bf Ayudantía 1}}
\centerline{\bf 10 de marzo de 2026}

\vspace{1cm}
\noindent\textbf{Problema 1.} Pruebe que dos rectas son paralelas si y solo si se cumple la 
propiedad de ángulos alternos internos. Esto es, dadas dos rectas paralelas y una recta secante a
ambas, entonces los ángulos alternos internos son iguales.

\vspace{5mm}
\noindent\textbf{Problema 2.} Sean $L_{1}$ y $L_{2}$ dos rectas paralelas. Se traza una recta 
$L_{3}$ secante a ambas rectas que intersecta a $L_{1}$ y $L_{2}$ en los puntos $A$ y $D$ 
respectivamente. Sea $L_{4}$ una recta que intersecta a $L_{3}$ en el punto $E$ que se encuentra 
entre ambas rectas paralelas y las intersecta en los puntos $C$ y $B$ en ese orden. Muestre que
\begin{equation*}
    \angle DBE+\angle CAE=\angle BEA \hhtext{y} \angle BDE+\angle ACE=\angle DEC
\end{equation*}

\vspace{5mm}
\noindent\textbf{Problema 3.} Considere el $\triangle ABC$ de modo que $\angle ABC=\angle ACB$. 
Sean $D$ y $E$ puntos sobre $\overline{BC}$ y $\overline{AC}$, respectivamente, tales que 
$\angle BAD=30^{\circ}$ y $\angle ADE=\angle AED$. Encuentre el valor del ángulo 
$\angle EDC$.

%\printbibliography % Quitar el comentado si quiero usar bibliografia

\end{document}
